\SourceLeaf{286}%
können, braucht man zwar Näherungswerte für die Längen der Linien, die orientierten Azimute und die Breiten der Endpunkte. Diese lassen sich nur durch eine vorläufige Berechnung der Dreiecke gewinnen. Das fällt aber überhaupt nicht ins Gewicht, da eine vorläufige Rechnung, die schon die Messungen Schritt für Schritt begleitet, ohnehin in vieler Hinsicht ratsam ist. Für die Zentrierung der exzentrisch gemessenen Winkel sowie für die Bestimmung des sphärischen oder sphäroidischen Exzesses der Winkelsumme jedes Dreiecks ist sie sogar notwendig. Für den ersten Zweck reicht bei der Geringfügigkeit jener Reduktionen stets schon eine sehr grobe Näherung aus, während die genaue Zentrierung bei etwas größerer Exzentrizität der Standpunkte bisweilen eine erheblich genauere Näherung erfordern kann. Ich habe die Vorschriften entwickelt, damit man, \textit{wenn} man jene Reduktionen berücksichtigen will, alles für ihre möglichst genaue Berechnung Nötige zur Hand hat oder, wenn man sie \textit{nicht} berücksichtigen will, leicht und genau erkennen kann, wie wenig man dadurch aufgibt. Im gesamten hannoverschen Dreiecksnetz sind die Reduktionen durchweg so außerordentlich gering, dass ihre Berücksichtigung völlig überflüssig erscheint. Über die ganzen zwölf Breitengrade, für die ich die Hilfsmittel beifüge, bleiben sie noch unter den Bruchteilen einer Bogensekunde, auf die man sich bei den meisten Messungen in der Rechnung gewöhnlich beschränkt. Um das besonders deutlich zu machen, füge ich hier noch die Zahlenrechnung für zwei Beispiele bei.

Im hannoverschen Dreiecksnetz treten die größten Reduktionen bei den Richtungen der Seiten des Dreiecks Brocken-Hohehagen-Inselsberg auf. Dieses Dreieck ist zugleich das größte und das vom Normalparallelkreis am weitesten entfernte. Bei allen übrigen Dreiecksseiten überschreiten die Reduktionen nirgends zwei Tausendstel einer Sekunde, und die meisten erreichen nicht einmal den Wert \(0''001\).

\Needspace{7\baselineskip}
Für diese Punkte gilt
\begin{center}
\begin{tabular}{l|c|c|r}
&\multicolumn{2}{c|}{Breite}&\\
&auf dem Ellipsoid&auf der Kugel&\(k\)\\\hline
Brocken&\(51^0\quad48'\quad2''\)&\(51^0\quad46'\quad3''\)&\(0''164\)\\
Hohehagen&\(51\quad28\quad31\)&\(51\quad26\quad35\)&\(0{,}303\)\\
Inselsberg&\(50\quad51\quad9\)&\(50\quad49\quad16\)&\(0{,}687\)
\end{tabular}
\end{center}
